\documentclass[a4paper,10pt]{article}
\usepackage[utf8]{inputenc}
\usepackage[T1]{fontenc}
\usepackage[sc]{mathpazo}
\usepackage[scaled=0.88]{helvet}
\usepackage{amsmath,amssymb,mathrsfs,stmaryrd}
\usepackage{microtype}
\usepackage[a4paper,top=1.5cm,bottom=1.5cm,left=1.4cm,right=1.4cm,headheight=12pt,headsep=0.4cm,footskip=0.7cm]{geometry}
\usepackage{xcolor}
\definecolor{Ink}{HTML}{1C2533}\definecolor{Bleu}{HTML}{1F3A5F}\definecolor{BleuMoyen}{HTML}{4A6A8F}
\definecolor{BleuClair}{HTML}{9DB2C9}\definecolor{Fond}{HTML}{F2F4F7}\definecolor{Ocre}{HTML}{B07A26}
\definecolor{Brique}{HTML}{A23B2A}\definecolor{Gris}{HTML}{5E6672}\definecolor{Filet}{HTML}{C5CCD6}
\usepackage{multicol,enumitem,booktabs,array,needspace}
\usepackage[most]{tcolorbox}
\usepackage{tikz,pgfplots}
\pgfplotsset{compat=1.17,trig format plots=rad}
\usetikzlibrary{arrows.meta}
\usepackage{fancyhdr}
\usepackage[colorlinks=true,linkcolor=Bleu,bookmarksopen=true,bookmarksopenlevel=0]{hyperref}

\setlength{\parindent}{0pt}\setlength{\parskip}{0.25em}
\setlength{\columnsep}{0.55cm}\setlength{\columnseprule}{0.3pt}\renewcommand{\columnseprulecolor}{\color{Filet}}
\setlist{itemsep=0.05em,topsep=0.15em,parsep=0pt,leftmargin=1.1em}
\setlist[itemize]{label=\textcolor{BleuMoyen}{\scriptsize\textbullet}}
\setlist[enumerate]{label=\textsf{\scriptsize\color{Bleu}\arabic*.},leftmargin=1.3em}
\linespread{1.02}

\newcommand{\K}{\mathbb{K}}\newcommand{\R}{\mathbb{R}}\newcommand{\C}{\mathbb{C}}\newcommand{\Z}{\mathbb{Z}}
\newcommand{\Mn}[1]{\mathcal{M}_{n}(#1)}\newcommand{\rg}{\operatorname{rg}}\newcommand{\vect}{\operatorname{Vect}}
\newcommand{\Mat}{\operatorname{Mat}}\newcommand{\Id}{\operatorname{Id}}\newcommand{\com}{\operatorname{com}}
\newcommand{\vv}[1]{\overrightarrow{#1}}\newcommand{\ve}[1]{\vec{#1}}
\newcommand{\dt}[1]{\left|\begin{matrix}#1\end{matrix}\right|}
\newcommand{\HP}{{\sffamily\bfseries\scriptsize\textcolor{Brique}{[HP]}}}
\newcommand{\livre}[1]{{\sffamily\scriptsize\color{Gris}[#1]}}

% ---------- titres ----------
\newcommand{\chapitre}[2]{\clearpage\phantomsection\pdfbookmark[0]{#1 --- #2}{chap#1}\markboth{#1 --- #2}{}%
  \noindent\colorbox{Bleu}{\parbox{\dimexpr\linewidth-2\fboxsep}{\color{white}\sffamily\bfseries\Large\strut #1\quad
  {\normalsize\mdseries #2}}}\par\vspace{0.4em}}
\newcommand{\notion}[1]{\par\needspace{9\baselineskip}\vspace{0.5em}\phantomsection\pdfbookmark[1]{#1}{\detokenize{#1}}\noindent{\sffamily\bfseries\color{Bleu}\large #1}\par\nobreak
  \noindent{\color{Ocre}\rule{\linewidth}{0.8pt}}\par\nobreak\vspace{0.1em}}
\newcommand{\sous}[1]{\par\needspace{6\baselineskip}\vspace{0.25em}\noindent{\sffamily\bfseries\small\color{BleuMoyen}#1}\par\nobreak}

% ---------- encadrés ----------
\tcbset{base/.style={enhanced,arc=0pt,boxrule=0pt,left=4pt,right=3pt,top=2pt,bottom=2pt,before skip=0.35em,after skip=0.35em,
  fontupper=\small}}
\newtcolorbox{DEF}[1][]{base,colback=white,borderline west={1.4pt}{0pt}{Bleu},
  before upper={{\sffamily\bfseries\scriptsize\color{Bleu}DÉFINITION\ifx\relax#1\relax\else\ --- #1\fi}\par}}
\newtcolorbox{PROP}[1][]{base,colback=Fond,borderline west={2pt}{0pt}{Bleu},
  before upper={{\sffamily\bfseries\scriptsize\color{Bleu}\ifx\relax#1\relax PROPRIÉTÉ\else #1\fi}\par}}
\newtcolorbox{FORM}[1][]{base,colback=white,colframe=BleuMoyen,boxrule=0.5pt,
  before upper={{\sffamily\bfseries\scriptsize\color{BleuMoyen}\ifx\relax#1\relax FORMULES\else #1\fi}\par}}
\newtcolorbox{METH}[1]{base,colback=white,borderline north={0.8pt}{0pt}{Bleu},borderline south={0.8pt}{0pt}{Bleu},
  before upper={{\sffamily\bfseries\scriptsize\colorbox{Bleu}{\color{white}MÉTHODE}\ \color{Bleu}#1}\par\smallskip}}
\newtcolorbox{PIEGE}{base,colback=Brique!5,borderline west={1.4pt}{0pt}{Brique},
  before upper={{\sffamily\bfseries\scriptsize\color{Brique}PIÈGES}\par}}
\newtcolorbox{ASTU}{base,colback=Ocre!7,borderline west={1.4pt}{0pt}{Ocre},
  before upper={{\sffamily\bfseries\scriptsize\color{Ocre}ASTUCES}\par}}
\newtcolorbox{HPBOX}{base,colback=white,borderline west={1.4pt}{0pt}{Brique},
  before upper={{\sffamily\bfseries\scriptsize\color{Brique}HORS PROGRAMME}\par}}

\pagestyle{fancy}\fancyhf{}
\fancyhead[L]{\scriptsize\sffamily\color{Gris}\leftmark}
\fancyhead[R]{\scriptsize\sffamily\scshape\color{Gris}Le Collé PT --- fiche Séries numériques}
\fancyfoot[C]{\scriptsize\sffamily\thepage}
\renewcommand{\headrulewidth}{0.3pt}

\pgfplotsset{mini/.style={axis lines=middle,axis line style={Filet},xtick=\empty,ytick=\empty,width=3.2cm,height=3.2cm,
  xmin=-1.15,xmax=1.15,ymin=-1.15,ymax=1.15,title style={font=\scriptsize\sffamily,yshift=-0.15cm},samples=80}}


\newcommand{\Ker}{\operatorname{Ker}}\renewcommand{\Im}{\operatorname{Im}}\newcommand{\tr}{\operatorname{tr}}\newcommand{\Sp}{\operatorname{Sp}}
\newcommand{\GLn}[1]{\mathrm{GL}_n(#1)}
\hypersetup{pdftitle={Le Collé — fiche de référence : séries numériques}}
\begin{document}
\thispagestyle{empty}
\noindent\colorbox{Bleu}{\parbox{\dimexpr\linewidth-2\fboxsep}{\vspace{0.4em}
{\color{white}\sffamily\scriptsize\textls[200]{LE COLLÉ\ \ ·\ \ MATHÉMATIQUES PT\ \ ·\ \ CHAPITRE 7}}\par\smallskip
{\color{white}\sffamily\bfseries\LARGE Fiche de référence --- Séries numériques}\par\smallskip
{\color{white!85!Bleu}\sffamily suites, séries de référence, termes positifs, convergence absolue, séries alternées, produit de Cauchy, résultats du TD}\vspace{0.4em}}}

\vspace{0.6em}
{\small Le livret sert à apprendre, la présentation à comprendre, cette fiche à travailler et réviser. Mêmes numéros de parties et
mêmes notations que la présentation et le livret ; \livre{th.~1.22} renvoie au livret, \livre{TD~8} à l'exercice~8 du TD complet.
$S_n=\sum_{k=0}^nu_k$ ; en cas de convergence, $S=\sum_{n=0}^{+\infty}u_n$ et $R_n=S-S_n$.}

\vspace{0.4em}
{\small\sffamily\noindent\begin{tabular}{@{}p{0.6\linewidth}@{}}
1.1 Rappels sur les suites \dotfill\pageref{n11}\\
1.2 Séries : définitions, premiers exemples \dotfill\pageref{n12}\\
1.3 Séries à termes positifs \dotfill\pageref{n13}\\
1.4 Convergence absolue, séries alternées \dotfill\pageref{n14}\\
1.5 Produit de Cauchy \dotfill\pageref{n15}\\
1.6 Quelle méthode ? Pièges ; résultats du TD \dotfill\pageref{n16}\\
\end{tabular}}

\vspace{0.6em}
\pdfbookmark[0]{Chapitre 7 --- Séries numériques}{chap7}\markboth{Chapitre 7 --- Séries numériques}{}
\begin{multicols}{2}\raggedcolumns
% ======================================================================
\notion{1.1 Rappels sur les suites}\label{n11}

\begin{PROP}[EXISTENCE \livre{th.~1.4, 1.5}]
Croissante majorée : converge. Croissante non majorée : $\to+\infty$.\\
Adjacentes ($a_n\nearrow$, $b_n\searrow$, $b_n-a_n\to0$) : même limite $\ell$, $a_n\leq\ell\leq b_n$.
\end{PROP}

\begin{DEF}[COMPARAISON \livre{déf.~1.6}]
$u_n\sim v_n$ : $\frac{u_n}{v_n}\to1$ ;\ $o$ : $\to0$ ;\ $O$ : borné.\quad $u_n\sim v_n\iff u_n=v_n+o(v_n)$.
\end{DEF}

\begin{FORM}[DL EN $0$]
$e^x=1+x+\frac{x^2}2+o(x^2)$,\quad $\ln(1+x)=x-\frac{x^2}2+\frac{x^3}3+o(x^3)$\\
$\cos x=1-\frac{x^2}2+o(x^3)$,\quad $\sin x=x-\frac{x^3}6+o(x^4)$\\
$(1+x)^\alpha=1+\alpha x+\frac{\alpha(\alpha-1)}2x^2+o(x^2)$,\quad $\frac1{1-x}=1+x+x^2+o(x^2)$
\end{FORM}

\begin{FORM}[CROISSANCES COMPARÉES \livre{th.~1.8}]
$(\ln n)^\alpha=o(n^\beta)$, $n^\alpha=o(a^n)$ ($a>1$), $a^n=o(n!)$, $n!=o(n^n)$.
\end{FORM}

\begin{PIEGE}
Pas de somme d'équivalents ; pas de composition ($e^{n+1}\not\sim e^n$). Puissance variable : $a_n^{b_n}=e^{b_n\ln a_n}$.
\end{PIEGE}

% ======================================================================
\notion{1.2 Séries : définitions, premiers exemples}\label{n12}

\begin{DEF}[\livre{déf.~1.10}]
$\sum u_n$ converge si $(S_n)$ converge. Nature inchangée si on modifie un nombre fini de termes. Reste $R_n\to0$.
\end{DEF}

\begin{PROP}[DIVERGENCE GROSSIÈRE \livre{th.~1.12}]
CV $\Rightarrow u_n\to0$. Réciproque fausse : $\sum\frac1n$.
\end{PROP}

\begin{FORM}[RÉFÉRENCES \livre{th.~1.14, 1.16, 1.17}]
$\sum(a_{n+1}-a_n)$ CV $\iff(a_n)$ CV ; somme $\lim a_n-a_0$.\\
$\sum q^n$ CV $\iff|q|<1$ ; $\sum_{n\geq0}q^n=\frac1{1-q}$.\\
$\sum\frac{x^n}{n!}=e^x$ ;\quad $\sum_{n\geq1}\frac1{n(n+1)}=1$.
\end{FORM}

% ======================================================================
\notion{1.3 Séries à termes positifs}\label{n13}

\begin{PROP}[\livre{th.~1.18 à 1.20}]
$u_n\geq0$ : CV $\iff(S_n)$ majorée.\\
$0\leq u_n\leq v_n$ : $\sum v_n$ CV $\Rightarrow\sum u_n$ CV.\\
$u_n\sim v_n$, \textbf{$v_n\geq0$} : même nature. $u_n=O(v_n)$ : $\sum v_n$ CV $\Rightarrow\sum u_n$ CV.
\end{PROP}

\begin{PROP}[SÉRIE-INTÉGRALE \livre{th.~1.22}]
$f$ continue, positive, \textbf{décroissante} : $f(k+1)\leq\int_k^{k+1}f\leq f(k)$ ;
$\sum f(n)$ et $\int^{+\infty}f$ ont même nature. Encadre aussi $S_n$ et $R_n$.
\end{PROP}

\begin{FORM}[RIEMANN \livre{th.~1.23}]
$\sum\frac1{n^\alpha}$ CV $\iff\alpha>1$.\quad $H_n\sim\ln n$ ;\ $H_n=\ln n+\gamma+o(1)$.\\
$\sum_{k>n}\frac1{k^\alpha}\sim\frac1{(\alpha-1)n^{\alpha-1}}$ \livre{TD~17}.
\end{FORM}

\begin{PROP}[D'ALEMBERT \livre{th.~1.25}]
$u_n>0$, $\frac{u_{n+1}}{u_n}\to\ell$ : $\ell<1$ CV ; $\ell>1$ DV ; $\ell=1$ ?
\end{PROP}

\begin{METH}{Termes positifs \livre{méth.~1.3}}
Équivalent $\to$ Riemann/géométrique. $n^\alpha u_n\to0$ ($\alpha>1$) : CV ; $nu_n\to+\infty$ : DV.
Factorielles : d'Alembert. $f$ décroissante : intégrale.
\end{METH}

% ======================================================================
\notion{1.4 Convergence absolue, séries alternées}\label{n14}

\begin{PROP}[\livre{th.~1.29, 1.30}]
$\sum|u_n|$ CV $\Rightarrow\sum u_n$ CV, $|\sum u_n|\leq\sum|u_n|$.\\
$u_n=O\left(\frac1{n^\alpha}\right)$, $\alpha>1$ : CV absolue.
\end{PROP}

\begin{PROP}[SÉRIES ALTERNÉES \livre{th.~1.32}]
$a_n\geq0$, \textbf{décroissante}, $\to0$ : $\sum(-1)^na_n$ CV ; $|R_n|\leq a_{n+1}$, signe du premier terme omis.
\end{PROP}

\begin{METH}{Signe quelconque \livre{méth.~1.4}}
Développer jusqu'à $O\left(\frac1{n^\alpha}\right)$, $\alpha>1$ ; conclure par linéarité.\\
$\frac{(-1)^n}{\sqrt n+(-1)^n}=\frac{(-1)^n}{\sqrt n}-\frac1n+O\left(\frac1{n^{3/2}}\right)$ : DV.
\end{METH}

\begin{PIEGE}
Équivalent pour un signe variable : interdit. Semi-convergentes : $\sum\frac{(-1)^n}n=-\ln2$, $\sum\frac{(-1)^n}{\sqrt n}$.
\end{PIEGE}

% ======================================================================
\notion{1.5 Produit de Cauchy}\label{n15}

\begin{PROP}[\livre{th.~1.36}]
$w_n=\sum_{p=0}^nu_pv_{n-p}$. Séries \textbf{absolument} CV : $\sum w_n=\sum u_n\cdot\sum v_n$.\\
$\sum(n+1)x^n=\frac1{(1-x)^2}$ ;\ $\sum nx^n=\frac x{(1-x)^2}$ ;\ $e^{a+b}=e^ae^b$.
\end{PROP}

% ======================================================================
\notion{1.6 Quelle méthode ? Pièges ; résultats du TD}\label{n16}

{\small\renewcommand{\arraystretch}{1.12}
\begin{tabular}{@{}>{\raggedright\arraybackslash}p{0.34\linewidth}>{\raggedright\arraybackslash}p{0.6\linewidth}@{}}
\hline
\textbf{Terme général} & \textbf{Outil}\\\hline
$u_n\not\to0$ & divergence grossière\\
fraction, racines & équivalent, Riemann\\
$\ln$, $\sin$\dots\ de $\to0$ & DL puis équivalent\\
$n!$, $a^n$, $n^n$ & d'Alembert\\
$f(n)$, $f\searrow$ & série-intégrale\\
$(-1)^n$ & absolue ; critère ; DL\\
$a_{n+1}-a_n$ & télescopique\\\hline
\end{tabular}}

{\small
\renewcommand{\arraystretch}{1.12}
\begin{tabular}{@{}>{\sffamily\scriptsize\bfseries\color{Bleu}}l>{\raggedright\arraybackslash}p{0.84\linewidth}@{}}
TD5 & $\sum\frac1{n(n+1)(n+2)}=\frac14$ ;\ $\sum_{n\geq2}\ln(1-\frac1{n^2})=-\ln2$.\\
TD6 & $\sum\frac1{n(n+2)}=\frac34$.\\
TD12 & $\sum kx^{k-1}=\frac1{(1-x)^2}$ ;\ $\sum(k-(-1)^k)3^{-k}=0$.\\
TD13 & $\sum\frac1{n(\ln n)^2}$ CV ;\ $\sum_{k\leq n}\frac1{k\ln k}\sim\ln\ln n$.\\
TD14 & $\sum\frac{(-1)^{n+1}}n=\ln2$ ; $10^3$ termes pour $10^{-3}$.\\
TD18 & $\ln(1+\frac{(-1)^n}{\sqrt n})$ : DV ;\ $\frac{(-1)^n}{n+(-1)^n}$ : CV.\\
TD19 & $\sum\frac1{n^\alpha\ln n}$ CV $\iff\alpha>1$.\\
TD20 & $\sum\frac{(-1)^n}{3n+1}=\frac{\ln2}3+\frac\pi{3\sqrt3}$.\\
TD23 & $R_n\sim\frac{(-1)^{n+1}}{2n}$ pour $\sum\frac{(-1)^k}k$.\\
TD25 & $u_{n+1}=\sin u_n$ : $u_n\sim\sqrt{3/n}$.
\end{tabular}}

\end{multicols}
\end{document}
